
\subsection*{Inverses After Restriction}

\begin{remark*}[Exposition]
Many functions are not invertible as originally stated. They may fail because
two inputs collide, because the codomain is too large, or both. Restricting
the domain can remove collisions; corestricting the codomain to the image
removes unused target elements.
\end{remark*}

\begin{theorembox}{Theorem (Injective Maps Are Bijections Onto Their Images)}
\begin{theorem}[Injective Maps Are Bijections Onto Their Images]\label{thm:injective-bijection-onto-image}
If $f:A\to B$ is injective, then the corestriction
$f_{\operatorname{im}}:A\to\im(f)$ is bijective. Hence it has an inverse
\[
f_{\operatorname{im}}^{-1}:\im(f)\to A.
\]
\hyperref[prf:injective-bijection-onto-image]{Go to proof.}
\end{theorem}
\end{theorembox}

\begin{remark*}[Standard quantified statement]
\[
  \operatorname{Injective}(f)\to \operatorname{Bijective}(f_{\operatorname{im}}).
\]
\end{remark*}

\begin{remark*}[Interpretation]
An injective function becomes bijective when its codomain is changed to its
image.
\end{remark*}

\begin{remark*}[Exposition]
The inverse-on-image is the function determined by
\[
\forall y\in\im(f)\,\exists! a\in A,\quad f(a)=y.
\]
Existence comes from $y\in\im(f)$; uniqueness comes from injectivity.
\end{remark*}

\begin{dependencies}
\begin{itemize}
  \item \hyperref[def:corestriction-to-image]{Corestriction to the Image}
  \item \hyperref[def:injective]{Injective Function}
  \item \hyperref[def:bijective]{Bijective Function}
  \item \hyperref[def:inverse]{Inverse Function}
\end{itemize}
\end{dependencies}

\begin{propositionbox}{Proposition (Restricted Map Is Onto Its Restricted Image)}
\begin{proposition}[Restricted Map Is Onto Its Restricted Image]\label{prop:restricted-map-onto-image}
Let $f:A\to B$ and let $C\subseteq A$. Then
\[
f|_C:C\to f(C)
\]
is surjective.
\hyperref[prf:restricted-map-onto-image]{Go to proof.}
\end{proposition}
\end{propositionbox}

\begin{remark*}[Standard quantified statement]
\[
  \forall y\in f(C),\ \exists c\in C,\quad f|_C(c)=y.
\]
\end{remark*}

\begin{remark*}[Predicate reading]
\[
  \operatorname{Surjective}(f|_C:C\to f(C)).
\]
\end{remark*}

\begin{remark*}[Interpretation]
Restricting the codomain to the restricted image makes the restricted map
surjective by definition.
\end{remark*}

\begin{dependencies}
\begin{itemize}
  \item \hyperref[def:restriction]{Restriction}
  \item \hyperref[def:image-set]{Image of a Set}
  \item \hyperref[def:surjective]{Surjective Function}
\end{itemize}
\end{dependencies}

\begin{propositionbox}{Proposition (Injective Restriction Gives a Local Inverse)}
\begin{proposition}[Injective Restriction Gives a Local Inverse]\label{prop:injective-restriction-local-inverse}
Let $f:A\to B$ and $C\subseteq A$. If $f|_C:C\to B$ is injective, then
\[
f|_C:C\to f(C)
\]
is bijective and therefore has an inverse $f(C)\to C$.
\hyperref[prf:injective-restriction-local-inverse]{Go to proof.}
\end{proposition}
\end{propositionbox}

\begin{remark*}[Standard quantified statement]
\[
  \operatorname{Injective}(f|_C)
  \to
  \operatorname{Bijective}(f|_C:C\to f(C)).
\]
\end{remark*}

\begin{remark*}[Interpretation]
Injectivity on a restricted domain plus corestriction to the restricted image
gives a local inverse.
\end{remark*}

\begin{dependencies}
\begin{itemize}
  \item \hyperref[prop:restricted-map-onto-image]{Restricted Map Is Onto Its Restricted Image}
  \item \hyperref[def:injective]{Injective Function}
  \item \hyperref[def:inverse]{Inverse Function}
\end{itemize}
\end{dependencies}

\subsection*{Failure Modes for Invertibility}

\begin{remark*}[Failure modes]
\begin{itemize}
\item[Exposition.]
To make $f:A\to B$ invertible, two independent failures must be repaired:
Domain restriction addresses collisions. Corestriction to the image addresses
missed targets.
\item[Collision failure.]
\[
  \exists a_1,a_2\in A,\quad a_1\neq a_2\wedge f(a_1)=f(a_2).
\]
This failure prevents injectivity.
\item[Missed-target failure.]
\[
  \exists b\in B,\quad \forall a\in A,\ f(a)\neq b.
\]
This failure prevents surjectivity.
\end{itemize}
\end{remark*}

\begin{remark*}[Negated quantified statement]
The negation of injectivity is not just ``not one-to-one''; it is the usable
predicate
\[
\exists a_1,a_2\in A,\quad a_1\neq a_2\ \land\ f(a_1)=f(a_2).
\]
The negation of surjectivity is
\[
\exists b\in B,\quad \forall a\in A,\ f(a)\neq b.
\]
\end{remark*}
